
\chapter{Glossary}

\begin{description}
    \item[Base] The number of distinct digits used in a positional numeral system.
    \item[Congruence] Two integers are congruent modulo $n$ if their difference is divisible by $n$.
    \item[Index (Exponent)] The power to which a base is raised; e.g.\ in $a^n$, $n$ is the index.
    \item[Logarithm] The inverse operation to exponentiation; $\log_b x = y \iff b^y = x$.
    \item[Set] A well-defined collection of distinct objects called elements.
    \item[Variation] A relationship between two quantities where one changes in proportion to the other.
    \item[Probability] A measure of how likely an event is to occur, expressed as a value between 0 and~1.
\end{description}

\textit{Expand this glossary as new terms are introduced throughout the book.}
